\section{Introducción}

Una transferencia de dinero parece una operación sencilla: restar de una
cuenta y sumar a otra. Casi toda la dificultad de un sistema de pagos está en
lo que ocurre cuando esa operación no sale según lo previsto. Qué pasa si el
proceso muere a la mitad, si dos peticiones llegan a la vez, si la respuesta
se pierde en la red y el cliente reintenta, o si alguien envía el
identificador de una cuenta ajena.

Este informe documenta el diseño, la construcción y la verificación de
PayCore, un prototipo educativo que aborda exactamente esas situaciones. El
proyecto nace como pieza de portafolio en Ciencia de la Computación, con un
criterio que ha gobernado todas las decisiones: \emph{poder explicar y
defender por qué cada cosa está hecha como está}.

\subsection{Enfoque}

Tres reglas de trabajo, fijadas antes de escribir la primera línea de código:

\begin{enumerate}[leftmargin=1.4em]
  \item \textbf{Investigar antes de programar.} Las decisiones sobre
    aislamiento, idempotencia o representación del dinero se tomaron después
    de leer la documentación oficial y la bibliografía correspondiente, no por
    costumbre. La investigación completa está en el repositorio y se resume en
    la sección~\ref{sec:investigacion}.
  \item \textbf{Verificar en los datos, no en las respuestas.} Un servicio que
    devolviera \codigo{201 Created} a todo sin escribir nada pasaría cualquier
    prueba que solo mire códigos HTTP. Por eso todas las aserciones sobre
    dinero consultan PostgreSQL directamente.
  \item \textbf{No añadir tecnología sin un requisito que la justifique.} No
    hay Redis, ni Kafka, ni Kubernetes, ni microservicios, ni modelos de
    aprendizaje automático. Para cada uno de ellos está escrita la condición
    concreta bajo la cual entraría.
\end{enumerate}

\subsection{Estructura del documento}

Las secciones \ref{sec:problema} a \ref{sec:investigacion} plantean el
problema y la investigación previa. Las secciones \ref{sec:tecnologias} a
\ref{sec:seguridad} describen las decisiones técnicas y su implementación. Las
secciones \ref{sec:pruebas} a \ref{sec:resultados} presentan la verificación y
las mediciones. El documento cierra con las limitaciones reconocidas, el
trabajo futuro y las conclusiones.
